\section{Related Work}
\label{sec:related}

\textbf{Interactive environments for coding.} 
Most coding benchmarks (e.g.\,SQL - Spider \cite{yu-etal-2018-spider}, KaggleDBQA \cite{lee-etal-2021-kaggledbqa}; Bash - NLC2CMD \cite{pmlr-v133-agarwal21b}, NL2Bash \cite{lin-etal-2018-nl2bash}; Python - HumanEval \cite{chen2021evaluating}, APPS \cite{hendrycksapps2021}, MBPP \cite{austin2021program}, CodeXGLUE \cite{lu2021codexglue}, CodeNet \cite{puri2021codenet}) frame the coding problem as a sequence transduction problem (from instruction to code), rather than an interactive decision making problem with an execution environment. 
Attempts have been made to simulate interaction by developing conversational, dialogue-style \cite{yu-etal-2019-sparc, yu2019cosql}, multi-step problem solving \cite{nijkamp2022codegen} datasets, which involve pre-annotated human-designed queries. The work closest to \benchmark{} has been recent explorations of Python Jupyter Notebooks as a natural choice for interactive coding \cite{huang-etal-2022-execution, Lai2022DS1000, yin2022natural}. However, task data and settings often constrain allowed actions to a closed domain of code and libraries~\cite{Lai2022DS1000, yin2022natural}, use evaluation procedures or metrics that may not generalize~\cite{huang-etal-2022-execution}, require human-in-the-loop participation (i.e. create task contexts, write problems, evaluate execution per task instance)~\cite{Lai2022DS1000}, or are Python-exclusive~\cite{huang-etal-2022-execution, Lai2022DS1000, yin2022natural, wang2023executionbased}.
\benchmark{} provides a more general purpose foundation defining interactive coding tasks that enables easy construction of diverse task settings, can have any programming language(s) as the action space, and has automatic, execution-based evaluation.

\textbf{Execution-based evaluation for coding.} Evaluation for NL-to-code generation models has recently shifted away from surface form similarity metrics (BLEU~\cite{Papineni2002BleuAM, Agashe2019JuICeAL}, ROUGE~\cite{lin-2004-rouge}, Exact Match) towards execution oriented ratings (unit tests~\cite{austin2021program, chen2021evaluating, huang-etal-2022-execution, Lai2022DS1000, hendrycksapps2021}, output matching~\cite{dong2016language, huang-etal-2022-execution, zhong2017seq2sql}).
The rigidity of surface form analysis overlooks code syntax features, ignores execution effect, or over-penalizes alternative solutions~\cite{zhou2023codebertscore},
On the contrary, execution-based assessment is a more thorough and comprehensive score of code functionality~\cite{hendrycksapps2021} and is a more natural fit for open-domain program usage that does not constrain code generation to a subset of the language space~\cite{wang2023executionbased}.
However, for newer benchmarks and datasets that put forth task definitions incorporating execution-based evaluation (APPS~\cite{hendrycksapps2021}, ExeDS~\cite{huang-etal-2022-execution},  ODEX~\cite{wang2023executionbased}), the fundamental code generation task (Context + Code $\rightarrow$ Execution $\rightarrow$ Score) is still devoid of interaction.
\benchmark{} combines execution-based evaluation with flexible task construction, enabling more diverse problem-solving paradigms within a unified coding task formulation.
\benchmark{}'s use of virtual containers as execution sandboxes protect against harmful actions and allow for advanced evaluation criteria beyond the aforementioned ones.

\textbf{Methods for interactive or execution-based coding.} The value of generative code models and interactive problem solving has motivated a recent proliferation of work to augment reasoning capabilities' of existing language models~\cite{yao2023react, shinn2023reflexion, wang2023planandsolve, yao2023tree, zhang2023planning, chen2023teaching} or propose new modeling techniques to tackle coding as a sequential decision making and reasoning tasks~\cite{bunel2018leveraging, Chen2018ExecutionGuidedNP, ellis2019write, alphacode, chen2022codet, le2022coderl}, of which evaluation is unit test based.
Approaches that leverage execution typically use re-ranking~\cite{zhang2022coder, ni2023lever, yin-neubig-2019-reranking, zeng2022nbest} or majority vote~\cite{Chen2018ExecutionGuidedNP, alphacode, shi2022natural} to decide on a final prediction.
Additional work also explores incorporating human-in-the-loop~\cite{chen2023improving, lahiri2022interactive}, compilers~\cite{wang2022compilable}, and text~\cite{wang2023leti} feedback.
A common thread among these contributions is that 1) the task setting can only provide the investigated form of feedback and 2) sought-after capabilities are exemplified by strong performance on favorably curated tasks and datasets, rendering comparisons across benchmarks tedious.
\benchmark{} has the potential to standardize the evaluation of these methods because 1) the interactive coding task is a conglomeration of many interesting interaction, reasoning, and decision-making challenges and 2) \benchmark{}'s task construction makes it possible to incorporate a wide variety of sources of feedback.

















